\documentclass[11pt,a4paper]{article}

\usepackage[utf8]{inputenc}
\usepackage[T1]{fontenc}
\usepackage[french]{babel}
\usepackage{lmodern}
\usepackage{microtype}
\usepackage[margin=2.2cm]{geometry}
\usepackage{parskip}
\usepackage{xurl}
\usepackage[hidelinks]{hyperref}
\usepackage{fancyhdr}
\usepackage{cite}
\setlength{\emergencystretch}{1.5em}
\setlength{\headheight}{14pt}
\clubpenalty=10000
\widowpenalty=10000
\brokenpenalty=10000
\pagestyle{fancy}
\fancyhf{}
\fancyhead[L]{\small Carnet de recherche}
\fancyhead[R]{\small Maxime DRÉAN}
\fancyfoot[C]{\thepage}
\newcommand{\weblink}[2]{\href{#1}{#2}}
\newcommand{\doi}[1]{\href{https://doi.org/#1}{doi:#1}}

\title{Carnet de recherche\\[0.3em]
\large Reconnaissance acoustique des oiseaux\\par apprentissage profond}
\author{Maxime DRÉAN\\\small Master 2 -- Informatique 2349}
\date{Septembre -- octobre 2026}

\begin{document}
\maketitle
\thispagestyle{fancy}

\section{Lectures et recherches}

Le sujet de départ porte sur l'identification automatique des oiseaux dans
des enregistrements naturels. J'ai retenu trois difficultés : le bruit ambiant,
les vocalisations simultanées et le changement de région. La question de veille
est donc : dans quelle mesure les modèles d'apprentissage profond permettent-ils
de détecter et d'identifier les espèces malgré ces conditions ? Le périmètre
comprend les chants, les cris, les modèles et les données d'évaluation ; il
laisse de côté la localisation spatiale et l'estimation d'abondance.

La revue de Stowell a servi à situer les méthodes et les difficultés de la
bioacoustique automatique\cite{stowell2022}. BirdNET a fourni un premier
exemple de système utilisable, et BirdSet un cadre pour comparer des modèles
sur des données aviaires\cite{kahl2021,birdset2025}. Xeno-canto, Macaulay
Library et BirdCLEF ont été repérés comme ressources possibles. Xeno-canto
a été retenu pour les premiers essais ; BirdSet a ensuite fourni les données
annotées nécessaires à une évaluation plus précise.

Les mots-clés ont suivi cette progression : \emph{bird sound recognition},
\emph{BirdNET}, \emph{bioacoustics deep learning}, puis \emph{noise robustness},
\emph{overlapping vocalizations} et \emph{soundscapes}. Le questionnement
s'est ensuite étendu au transfert avec \emph{audio embeddings},
\emph{self-supervised learning} et \emph{domain generalization}.
Les lectures ont été prolongées au fil des essais, notamment lorsque les
résultats ont soulevé la question du seuil et du changement de jeu de données.

Les travaux de Ghani et al. ont orienté la comparaison vers les représentations
réutilisables et les paramètres réellement ajustés\cite{ghani2023,ghani2025}.
AVES, Bird-MAE et Perch 2.0 ont été retenus pour comparer des objectifs de
préentraînement et des modes d'adaptation différents
\cite{aves2023,birdmae2025,perch2025}. Les évaluations de Sethi et de Funosas
ont complété cette lecture par les conditions d'usage sur le terrain
\cite{sethi2024,funosas2026}. Les dépôts et fiches de modèles ont servi à
vérifier les formats, les outils disponibles et les licences. La synthèse
retient surtout les travaux qui éclairent ces choix, plutôt qu'une succession
de descriptions de modèles.

\section{Première expérimentation}

J'ai choisi BirdNET pour disposer d'un modèle préentraîné et tester directement
l'effet du bruit, des mélanges et du seuil. Les essais reposent d'abord sur
3\,591 enregistrements Xeno-canto. L'espèce principale annoncée fournit une
cible à rechercher ; les autres sorties sont conservées pour explorer le
comportement du modèle.

L'expérience de bruit porte sur 2\,937 enregistrements dont la cible est
retrouvée au départ. Le segment où son score est maximal reçoit un bruit
gaussien à 20, 10, 5 et 0~dB, avec cinq réalisations par niveau. Les mélanges
portent sur 1\,416 paires, aux coefficients 50/50, 75/25 et 90/10. Un autre
essai fait varier le seuil pour observer quelles sorties sont conservées.

La collecte et l'analyse ont été organisées par lots pour reprendre le
traitement sans recommencer les calculs terminés. Les difficultés d'exécution
ont concerné notamment CUDA/cuDNN et les entrées ONNX de durée variable.
Un signal non silencieux a servi à vérifier le fonctionnement de l'inférence.
Le notebook a été conservé pour réunir les tableaux, les figures et les
observations ; les calculs ont été regroupés dans l'application Python.

Les premières sorties montrent une diminution des scores sous bruit et une
récupération conjointe plus difficile dans les mélanges. Elles font surtout
apparaître un problème de protocole : sélectionner une cible déjà reconnue
et son meilleur segment favorise le point de départ. L'étiquette principale
Xeno-canto ne permet pas non plus de classer toutes les autres prédictions
en présences réelles ou en fausses alertes. J'ai donc complété l'exploration
par des données annotées et des témoins acoustiques.

\section{Amélioration de l'expérimentation}

Deux sous-ensembles de BirdSet ont été retenus : NBP, issu de NIPS4BPlus,
et HSN, issu de \emph{High Sierras}. NBP apporte 539 extraits annotés ;
HSN apporte 600 segments provenant de 100 sources, dont 333 segments sans
annotation d'oiseau. Ces derniers permettent d'observer les alertes en
l'absence de présence de référence. BirdNET 2.4 et 3.0 sont comparés sur
les mêmes extraits, sans filtrage géographique.

L'harmonisation des espèces a demandé de rapprocher les annotations des
vocabulaires des deux versions. Une espèce absente de 2.4 est comptée comme
omission ; un second calcul porte sur les classes communes. La comparaison
à seuil identique est complétée par un réglage propre à chaque couple
modèle/jeu : 25\,\% des sources servent au choix du seuil, les autres à
l'évaluation. Les segments d'un même enregistrement HSN restent groupés.
Les intervalles de confiance sont calculés par bootstrap apparié par source.

Le protocole acoustique a également été revu. Les 40 extraits positifs de
chaque jeu sont choisis à partir des annotations, en gardant les omissions
initiales. Des fonds enregistrés complètent le bruit gaussien. Pour les
mélanges, les niveaux efficaces sont égalisés et les deux signaux isolés
sont évalués avec les mêmes gains que dans le mélange. Les 20 paires NBP
et 19 paires HSN permettent ainsi de comparer la récupération conjointe
avant et après superposition.

Cette seconde étape a changé l'interprétation des essais. BirdNET 3.0
obtient un meilleur F1 sur NBP ; sur HSN, le seuil modifie fortement le
compromis entre omissions et fausses alertes. Après réglage séparé, aucun
avantage n'est établi sur HSN dans la partition retenue. Les témoins
confirment une baisse sous les perturbations étudiées, y compris lorsque
les extraits sont choisis par annotation. La variabilité géographique
reste discutée à partir des publications : les différences entre NBP
et HSN combinent aussi des espèces et des ambiances différentes.

\section{Rédaction du rapport final}

Le rapport final est organisé autour de la question de veille,
de la comparaison des méthodes et des résultats qui aident à choisir un
outil. L'état de l'art distingue la reconnaissance directe du transfert de
représentations, puis relie ces méthodes aux annotations, au masquage et au
changement de domaine. Les descriptions détaillées du code sont laissées
dans le rendu de l'expérimentation.

Pour la partie expérimentale, j'ai retenu la comparaison des versions et
des seuils, le bruit ajouté et les mélanges avec témoins. Le corpus Xeno-canto
explique le point de départ et l'évolution du protocole. Les tableaux et
figures mettent en avant les résultats utiles à l'argumentation, avec les
effectifs et les intervalles nécessaires à leur lecture.

La discussion confronte les observations aux travaux sur le transfert et
aux évaluations de terrain. Elle distingue les résultats mesurés des
explications possibles, par exemple pour l'écart entre bruit gaussien et
fonds enregistrés. La conclusion revient au choix conjoint du modèle,
des données et du seuil. Les prolongements retenus sont un test multisite
sur des espèces communes et une comparaison d'adaptations avec le même
budget d'annotations.

\bibliographystyle{unsrt}
{\small\bibliography{references}}

\end{document}
